% Overleaf-Datei, Spiegel in docs/overleaf/. Nur zusammen mit Overleaf aendern.
% Inhalt: Section 4.3 "Multireference geometries" (sec:mr-whynotmr, schliesst
% das OMol25-Kapitel ab) und Chapter 5 "Conclusion & Outlook" (ch:conclusion).
% Entwuerfe: docs/section_mr_outofscope.tex, docs/discussion_delta.md,
% docs/discussion_omol25.md.

\section{Multireference geometries: attempted, out of scope}
\label{sec:mr-whynotmr}

Multireference geometries and energies were attempted for these reactions
because they would be the gold standard for reactions with high
multireference character. They were placed out of scope. Problems arose,
and they are not needed for the question of this thesis: the models were
trained on \mbox{$\omega$B97M-V/def2-TZVPD} and are trying to find the
transition states of that surface, so the reference stays at that level.
A multireference reference
would test a different question: whether the models agree with a
surface they were never trained on. Where they did, it would be by
accident. No multireference barrier enters this work.

\paragraph{The overall problem: active space selection.} Every
multireference method needs an active space, and the result depends on
it. autoCAS~\cite{stein2019autocas} selects the space from a DMRG
(density matrix renormalisation group) calculation over all valence orbitals, which was not practical for
molecules of this size. Smaller molecules in this part of the test set
would have avoided this. AVAS (atomic valence active space)~\cite{sayfutyarova2017} and selection by
hand require experience that could not be acquired within the time of
this work and was not part of the thesis plan either. NEVPT2, the
standard protocol for energies on unrestricted geometries, builds on a
CASSCF wavefunction and therefore needs the same active space, so it
was not carried out either.

\paragraph{The additional problem for obtaining multireference-optimised
geometries: active space consistency.} Starting from the DFT transition state, AVAS selected a
space, CASSCF optimised a transition state with that active space, and
when AVAS was run again at the optimised geometry it returned a
different space. The active space that
fits a geometry changes as the geometry moves, so the surface the
optimisation ran on is not the surface that fits the geometry it ended
at. A possible solution is an iteration:

\begin{enumerate}
  \item select the active space, with autoCAS where affordable;
  \item optimise the transition state at the CASPT2 or NEVPT2 level;
  \item re-select the active space at the optimised geometry;
  \item repeat until the selected space is the one the geometry was
    optimised with;
  \item add the next orbitals to the converged space and check that the
    result does not change.
\end{enumerate}

The re-selection in the optimised orbital basis is what the autoCAS
authors recommend for a single structure~\cite{stein2019autocas}; the
iteration above applies it to a moving geometry.

A significant part of the time of this work went into this attempt.
Benchmark sets with transition-state geometries optimised at a
multireference level are rare; the standard one holds eleven
pericyclic hydrocarbon reactions with CASSCF transition structures and
CASPT2 energies~\cite{guner2003}, and nothing comparable exists for
reactions of the size and diversity of Transition1x. The iteration
above might be one way such a set could be attempted. For the question
of this thesis it was not needed: the models are trained on the \mbox{$\omega$B97M-V/def2-TZVPD}
surface and are tested against it, and that reference stands on its
own.

\chapter{Conclusion \& Outlook}
\label{ch:conclusion}

\section{Research question 1: raising the level of theory with a subset of the labels}

We return to Chapter~\ref{ch:delta} and RQ1, whether a model can be
lifted to a higher level of theory by relabelling only a subset of its
training data. The answer is yes, certainly as long as the spin
formalism is not changed with the labels, and Allen et
al.~\cite{allen2026} and DeePEST-OS~\cite{ren2026} have shown the
principle in related settings. With under 1\,\% of the dataset relabelled, the following
was achieved:

\begin{itemize}
  \item the force error drops on every reaction, to about a third in
    every tier;
  \item the energies improve in every tier;
  \item the geometries stay as good as MACE's: the transition states
    found with MACE+$\Delta$ lie as close to the reference as those
    found with MACE (mean RMSD 0.054 against 0.056\,\AA,
    Figure~\ref{fig:rmsd-parity}). No improvement was to be made
    here, since the transition states of
    \mbox{$\omega$B97X/6-31G(d)} and \mbox{$\omega$B97M-V/def2-TZVP}
    lie a median of 0.007\,\AA{} apart, below the geometry error
    of the base model.
\end{itemize}

The barriers depend on the multireference character. On the low-MR
reactions the barrier error drops to a third, in single points and in
searches the model drives itself. At mid and high multireference
character the correction overshoots, and the overshoot is mostly the
base model's error, which grows with multireference character and
passes through the head unchanged. Chapter~\ref{ch:mr} showed that
at mid MR the restricted reference is sound, while at high MR it is
unstable at seven of ten transition states. The errors at high MR are
still errors of the model: it was trained on restricted labels and
misses the restricted reference. Independently of that, at those seven
transition states the restricted reference is not the lowest DFT
solution, so a model that reproduced it exactly would still be wrong
there.

The correction is bounded by the base model. MACE alone misses its
own level by 42\,meV, that error enters MACE+$\Delta$ unchanged, and
the head cannot reduce it (Eq.~\eqref{eq:error-split}); the corrected
model ends at 56\,meV. The base model is weakest where the
multireference character is high: its barrier error against its own
level grows from $-91$\,meV at mid MR to $-192$\,meV at high MR. This
is not specific to MACE. The OMol25 models of Chapter~\ref{ch:mr},
which differ from it in architecture, training set and level of
theory, also degrade with multireference character, even at
closed-shell transition states, which suggests that the difficulty is
more fundamental than the choice of base model.

Four changes could have made the correction better.

\begin{itemize}
  \item \textbf{A better base.} Every improvement of MACE on its own
    level transfers one to one, because that term passes through the
    sum unchanged. At high multireference character this may meet a
    ceiling: as noted above, the OMol25 models show the same weakness,
    so a better base of the current kind may not remove it.
  \item \textbf{A trainable encoder.} The head reads frozen features,
    which is why the baseline passes through. Unfreezing the last MACE
    layer or the readout on the same 80\,000 points would let the model
    move its own surface towards the target instead of adding an
    offset. The risk is forgetting the cheap-level surface where the
    relabelled points are sparse.
  \item \textbf{An input that sees the gap change.} The head predicts
    the difference from scalar features alone. A feature that tracks
    multireference character, such as $N_\text{FOD}$ or the
    $\langle S^2 \rangle$ of the label, would let the head learn that
    the gap shrinks where it still over-corrects, at mid MR.
  \item \textbf{Labels on the right surface.} Where the restricted
    solution is unstable, the head learns a difference between two
    restricted energies. Relabelling those points unrestricted, with a
    stability analysis, would give it a difference that describes the
    unrestricted surface.
\end{itemize}

The limits of Chapter~\ref{ch:delta} are its 30 reactions, the restricted
treatment throughout, and def2-TZVP without the diffuse functions of
the target level.

The correction was not pushed further, because the principle had been
shown in related settings by others~\cite{allen2026, ren2026}; the open question
was where relabelling stops, which is what Chapter~\ref{ch:mr} took up.

\section{Research question 2: learning a surface from paths never relaxed on it}

RQ2 asked whether a model can learn a surface by relabelling paths
that were never relaxed on that surface. It does when only the level
of theory changes. It does not when the spin formalism changes as
well: the models then deliver structures that are not stationary
points of the new surface. OMol25 is
that case: it relabelled the whole Transition1x set unrestricted at
\mbox{$\omega$B97M-V/def2-TZVPD} without re-optimising a path.
Chapter~\ref{ch:mr} tested three of its models on 45 Transition1x
reactions and found:

\begin{itemize}
  \item at broken-symmetry transition states the models deliver
    structures that are not stationary points of the unrestricted
    surface, with DFT residual forces 2.5 times those of the
    closed-shell group, while their own forces report convergence;
  \item their force error there is two to three times larger, while
    the energies stay within chemical accuracy;
  \item the barrier error is a geometry error: the spread of the DFT
    barrier over the three model geometries is 32 times larger, 14
    times where all searches converged;
  \item the training data offer an explanation: on the unrestricted
    surface the stored transition states of the broken-symmetry
    reactions carry almost three times the residual force of the
    closed-shell ones, and the change of spin formalism, not the
    change of level, is responsible.
\end{itemize}

Part of this difference is multireference character itself, which
degrades the models' predictions at closed-shell transition states
too. Chapter~\ref{ch:delta} saw the same for MACE on its own level,
where no change of surface is involved. But when only reactions of the
same multireference character are
compared, the broken-symmetry structures are still worse, by a factor
of about 1.6. That part is the effect of broken symmetry itself. The
training data offer an explanation for it, and the shape of the
unrestricted surface is a second candidate.

Figure~\ref{fig:pes} gives intuition for what happens when a
broken-symmetry transition-state region is relabelled on cheap
restricted paths in the training data: a schematic of the same
geometries on the restricted surface they were relaxed on and on the
unrestricted surface the labels were computed on. On the first the
stored path holds the transition state; on the second it passes
beside it.

No reaction path in OMol25 was relaxed at the level its labels use
(Section~\ref{sec:omol-discussion}), so a saddle of the unrestricted
surface at the OMol25 level does not exist anywhere in its training
data.

\paragraph{A second candidate: the shape of the unrestricted surface.}
The training data are not the only possible cause. Where the
broken-symmetry solution sets in, the curvature of the unrestricted
surface changes abruptly, and a smooth model may fit that region poorly
whatever geometries it is trained on. Inside the broken-symmetry group
the data do not settle the matter: the reactions whose training points
are furthest from stationarity do not show the largest errors, and the
force error falls to the closed-shell level only above a breaking depth
of about 1.3~eV (Section~\ref{sec:omol-discussion}). Both causes may
contribute.

\paragraph{Limitations.} The sample is 45 reactions of small CHNO
molecules selected for multireference character; the findings concern
this class, not OMol25 as a whole. Of the four OMol25 models
benchmarked by Marks et al., three were tested here; MACE-OMol25, the
one with the highest success rate there, was not. This work contains no transition
state relaxed on the broken-symmetry surface, no multireference
reference (Section~\ref{sec:mr-whynotmr}), and no test of the models
after fine-tuning on such saddles, the test named below that would
separate the two candidate causes.

\section{Implications and outlook}

\paragraph{Why the restricted default persists.} We see three possible reasons why
the field seems to overlook the unrestricted treatment when it builds
datasets and benchmarks:

\begin{itemize}
  \item it is the default: a quantum chemistry program runs a singlet
    restricted unless told otherwise, and nothing in a high-throughput
    workflow tells it otherwise;
  \item the alternative is expensive and fragile;
  \item nothing in the workflow reveals the problem.
\end{itemize}

The first needs no explanation. Regarding the second point: an
unrestricted path needs a stability analysis or a carried-over guess
at every image, the SCF solution changes character along the path, and
images fall back to the restricted solution or jump between broken
ones. The attempts of this work at broken-symmetry transition states
met each of these and produced no reference that could be trusted. The
third is what makes the default persist: a restricted saddle is a
valid stationary point, the band converges, the frequency calculation
shows one imaginary mode, the model trains, and a benchmark built the
same way scores it correct. The default fails only against a check
that nobody in the chain runs.

\paragraph{The benchmarks are run restricted too.} The benchmarks the models
are tested against are run restricted, even where they contain
reactions that should be treated unrestricted. The reference set of
\'Asgeirsson et al.~\cite{asgeirsson2021} holds 27 such reactions
among 121, which Hait et al. found when they re-optimised it
unrestricted and then removed~\cite{hait2025}. The benchmark of Marks
et al.~\cite{marks2026} holds one among 25, optimised restricted like
the rest (Appendix~\ref{app:marks}). If the reference saddle is
restricted, a model that finds the restricted saddle passes, even
where the restricted solution is unstable.

Two consequences follow: one for benchmarks, and one for anyone who
runs a transition-state search. The question of the cause is left to
future work.

\paragraph{Benchmarks: a test set at broken-symmetry transition states.}
No benchmark exists that tests models systematically at
broken-symmetry transition states. We see this as a gap in the
literature: the reported success
rates were measured on restricted references, and it is easy to assume
that they hold for all transition-state searches.
This work shows that they do not, and closes the gap indirectly: the
residual force at the structure a model delivers needs no reference
saddle. A set of saddles optimised on the broken-symmetry
surface and confirmed by a stability analysis would close it directly
and give future models a benchmark to be tested against.

\paragraph{Practice: a stability analysis before trusting a saddle.} One unrestricted single point with a
stability analysis tells whether the restricted solution is the ground
state at a geometry, at the cost of one extra SCF. We recommend it at
every transition state that is trusted further, as a training point,
a reference, or a starting point, and whose multireference character
cannot be ruled out.

\paragraph{Future work: separating the two causes.} Whether the
training data or the shape of the unrestricted surface causes the
larger errors at broken-symmetry transition states is left open by this
work. Two steps would settle it. The first is more data: with 18
broken-symmetry reactions the comparisons inside the group are not
conclusive, and a larger set that covers the range of breaking depths
evenly would show whether the error follows the distance of the
training points from stationarity or the closeness to the onset of
symmetry breaking. The second is the direct test: fine-tuning the
models on paths relaxed on the unrestricted surface, one unrestricted
CI-NEB per affected reaction, with a stability analysis at the saddle
to confirm that the path was relaxed on the broken-symmetry surface and
not on the restricted one. If the training data are the cause, the
errors should disappear, and such paths are what a dataset like OMol25
needs. If the shape of the surface is the cause, they should remain.

In any case, the saddles optimised on the broken-symmetry surface would
also work as a benchmark for future models to be tested against.
